\documentclass[11pt]{article}
\usepackage[T1]{fontenc}
\usepackage{lmodern,amsmath,amssymb,amsthm,mathtools,array,longtable}
\usepackage[margin=1.05in]{geometry}
\usepackage[hidelinks]{hyperref}
\newtheorem{theorem}{Theorem}[section]
\newtheorem{proposition}[theorem]{Proposition}
\newtheorem{lemma}[theorem]{Lemma}
\newtheorem{corollary}[theorem]{Corollary}
\newcommand{\R}{\mathbb R}\newcommand{\Z}{\mathbb Z}
\newcommand{\E}{\mathbb E}\newcommand{\Var}{\operatorname{Var}}
\newcommand{\code}[1]{{\small\nolinkurl{#1}}}
\title{Energy statistics, signed transfer and Gaussian dynamics\protect\\in an exact finite triad}
\author{Project Stardust}
\date{Reviewed diagnostic draft, October 30, 2026}
\begin{document}
\maketitle
\begin{abstract}
We reconstruct a finite real solenoidal Fourier family and its exact
weighted sector-transfer ratio. Randomizing one signed amplitude by a
standard Gaussian gives a nonzero odd transfer observable. Every even
observation has zero conditional mean transfer, although weighted energy
determines its square and hence its strictly positive conditional variance.
Positive bounded changes of probability density preserve the entire joint
law of energy and concentration while reversing mean signed transfer.
Complete proofs distinguish equality of observed laws from equality of
latent laws. We also derive the biased conditional means and a two-law prediction-risk identity as explicit classical consequences. The actual retained dynamics has a nonlinear Gaussian-preserving global flow, a nonzero divergence-free quadratic Gaussian drift orthogonal to every original closed gradient, and coordinate dephasing despite a conserved observable and unitary $L^2$ action. The viscous counterpart dissipates every trajectory while admitting
unbounded pair amplification over all initial amplitudes, and loses
the inviscid Gaussian law. The value
is an exact diagnostic and formalization companion: phase dependence and
symmetry-based information loss are classical, and priority for this
particular package remains unresolved. No turbulent evolution, invariant
Navier--Stokes measure or failure of a dynamically constrained closure is
asserted.
\end{abstract}

\section{What is observed and what is tested}
Spectral energy is quadratic in Fourier coefficients, whereas a nonlinear
transfer is cubic. Sign information can therefore disappear from energy
statistics. We make this statement exact on an explicitly specified family,
partition and probability law, including all normalizing denominators.
The signed parameter below changes relative signs of imaginary Fourier
coefficients; it is not an angle uniformly distributed on a circle.

The source calls $\mathcal E$ an energy. It is the weighted quantity
$\sum_k|k|^2|\widehat u(k)|^2$, proportional to enstrophy on this
solenoidal family. It is not the unweighted kinetic energy. Our central
observation is $S=(\mathcal E,\mathcal C)$, with the concentration
$\mathcal C$ defined below. Its entire probability law will be preserved,
not just one or two moments.

The first six sections evaluate finite formulas on static fields. The later dynamical sections specify the retained ODE separately. We define a transfer
diagnostic by its displayed convolution formula. We do not assume that
the six retained modes remain closed under Euler or Navier--Stokes
evolution. Moffatt's primary analysis~\cite{moffatt}, introduction and
Section 2, makes precisely this distinction between triad truncation and
exact fluid evolution. None of his evolution results is needed for the
finite identities here.

\section{A real periodic triad and its complete finite transfer}
Let $e_1=(1,0,0)$, $e_2=(0,1,0)$ and
$K=\{\pm e_1,\pm e_2,\pm(e_1+e_2)\}\subset\Z^3$.
For real $a,b,c$ define a real vector polarization $V$ by
\[
 V(e_1)=(0,a,0),\quad V(e_2)=(-b,0,0),\quad
 V(e_1+e_2)=(-c,c,0),\quad V(-k)=-V(k),
\]
Set $V=0$ off $K$. Use complex Fourier coefficients $iV(k)$ on
the period-$2\pi$ three-torus, so that
\[u(x)=\sum_{k\in K}iV(k)e^{ik\cdot x}.\] Pairing opposite modes gives the actual field
\begin{equation}\label{eq:velocity}
 u_{a,b,c}(x)=\bigl(2b\sin x_2+2c\sin(x_1+x_2),
           -2a\sin x_1-2c\sin(x_1+x_2),0\bigr).
\end{equation}
Thus $u$ is real, smooth and periodic. Each $k\cdot V(k)=0$;
equivalently its two displayed divergence terms are
$2c\cos(x_1+x_2)$ and $-2c\cos(x_1+x_2)$. This proves
incompressibility without a pressure assumption. The field vanishes
identically only if $a=b=c=0$: evaluate its components at
$(\pi/2,0,0)$ and $(0,\pi/2,0)$ to obtain first $c=0$ and then
$a=b=0$.

For $k\in K$ define
\begin{equation}\label{eq:weights-transfer}
 W_k=|k|^2|V(k)|^2,\qquad
 T_k=2|k|^2\sum_{\substack{p,q\in K\\p+q=k}}
       (q\cdot V(p))(V(k)\cdot V(q)).
\end{equation}
Because $V$ is zero off $K$, the sum agrees with the unrestricted
convolution at each retained output frequency: an outside input has a
zero vector factor. This statement does not assert vanishing at other
output frequencies or invariance of the finite support under evolution.

\begin{lemma}\label{lem:table}
In the order $e_1,e_2,e_1+e_2,-e_1,-e_2,-e_1-e_2$, the weights
and transfers are
\[
 (W_k)=(a^2,b^2,4c^2,a^2,b^2,4c^2),
\quad
 (T_k)=(2abc,-2abc,0,2abc,-2abc,0).
\]
In particular $\sum T_k=0$ and
$\mathcal E:=\sum W_k=2a^2+2b^2+8c^2$.
\end{lemma}
\begin{proof}
The weight formula follows from $|e_1|^2=|e_2|^2=1$ and
$|e_1+e_2|^2=2$. For output $e_1$, the only ordered pairs are
$(e_1+e_2,-e_2)$ and $(-e_2,e_1+e_2)$.
The first summand vanishes because
$V(e_1)\cdot V(-e_2)=0$; the second is $b(ac)=abc$.
The prefactor is two, giving $2abc$. For output $e_2$, the pairs
are $(e_1+e_2,-e_1)$ and $(-e_1,e_1+e_2)$.
The first has a zero vector inner product and the second is
$(-a)(bc)=-abc$, giving $-2abc$. For output $e_1+e_2$, the
pairs $(e_1,e_2)$ and $(e_2,e_1)$ contribute $abc$ and $-abc$.
Their sum is zero. Replacing $(k,p,q)$ by $(-k,-p,-q)$ leaves
each summand unchanged: both dot products contain two sign reversals.
This proves the three opposite-output values and completes every
possible pair. Summing the two rows gives the last assertions.
\end{proof}

\section{An actual directional partition and the ratio}
We specify the source's 32 labels rather than infer a tessellation from
the six sample directions. For a direction $(x,y,z)$, assign a quadrant
$Q\in\{0,1,2,3\}$ by
\[
 Q=\begin{cases}0&y\geq0,\ x\geq0,\\1&y\geq0,\ x<0,\\
 2&y<0,\ x<0,\\3&y<0,\ x\geq0.\end{cases}
\]
Assign a bin $B\in\{0,\ldots,7\}$ by the first satisfied test in
the following ordered list, and assign $B=7$ if none is satisfied:
\[\begin{gathered}
 4|y|<|x|,\quad2|y|<|x|,\quad4|y|<3|x|,\quad|y|<|x|,\\
 3|y|<4|x|,\quad|y|<2|x|,\quad|y|<4|x|.
\end{gathered}\]
The sector is $s=8Q+B$. This rule assigns each point of the unit
sphere exactly one label and is invariant under positive scaling.
Every label is attained: in quadrant zero take $(1,r,0)$ with
\[
 r=1/8,\ 3/8,\ 5/8,\ 7/8,\ 7/6,\ 3/2,\ 3,\ 5
\]
for the eight bins, change the signs to enter the other quadrants,
and normalize to unit length. Normalization preserves every strict
comparison and quadrant. Thus these are disjoint nonempty cells
covering the sphere; no smooth-boundary or equal-area claim is needed.
On the six frequencies above their labels are
\[
 (s(k))=(0,7,4,8,31,20).
\]
They are distinct, even though the general rule has 32 cells.

Put $E_j=\sum_{s(k)=j}W_k$, $F_j=\sum_{s(k)=j}T_k$ and, for
$\mathcal E>0$, define
\begin{align}\label{eq:diagnostic}
 p_j&=E_j/\mathcal E, &\mathcal C&=\sum_{j=0}^{31}p_j^2,\notag\\
 \Sigma&=\frac2{\mathcal E}\sum_{j=0}^{31}
             p_j\left(F_j-p_j\sum_kT_k\right),
 &R&=-\frac{\Sigma}{\mathcal C-1/32}.
\end{align}
This precise formula defines the diagnostic being tested. Its
denominator and its sign will be verified on actual amplitudes.

\begin{proposition}\label{prop:ratio}
For $(a,b,c)\ne0$ let $E=2a^2+2b^2+8c^2$ and
$A=a^4+b^4+16c^4$. Then
\begin{equation}\label{eq:ratio}
 \mathcal C=\frac{2A}{E^2}\geq\frac16>\frac1{32},\qquad
 \Sigma=\frac{8abc(a^2-b^2)}{E^2},\qquad
 R=\frac{-8abc(a^2-b^2)}{2A-E^2/32}.
\end{equation}
The denominator is strictly positive. Swapping $a,b$ preserves
$E,\mathcal C$ and reverses $R$. At $(2,1,1)$ and $(1,2,1)$ the
ratios are $-128/149$ and $128/149$. At $(2,2,1)$ the ratio is
zero although $T_{e_1}=8$.
\end{proposition}
\begin{proof}
The six displayed sector weights are the six entries of Lemma~\ref{lem:table}.
Their squared sum is $2A$, proving the concentration formula.
Cauchy--Schwarz for these six weights gives
$E^2\leq6(2A)$, hence the displayed lower bound. Because $E>0$,
$2A-E^2/32=E^2(\mathcal C-1/32)>0$.
The total transfer vanishes, while
\[
 \sum_jE_jF_j=2a^2(2abc)+2b^2(-2abc)=4abc(a^2-b^2).
\]
Substitution in~\eqref{eq:diagnostic} gives $\Sigma$ and then $R$.
The numerator is antisymmetric in $a,b$ and its denominator symmetric.
At $(2,1,1)$, $E=18$, $A=33$ and
$2A-E^2/32=447/8$, giving $-48/(447/8)=-128/149$.
At $(2,2,1)$, $E=24$, $\mathcal C=1/6$, and the imbalance
$abc(a^2-b^2)$ vanishes, while the table gives $T_{e_1}=8$.
Thus zero normalized rate does not mean that every individual transfer
vanishes. All assertions follow from the actual finite table.
\end{proof}

\section{Gaussian signed amplitudes and exact invisible fluctuations}
Now fix $a=2,b=1$ and let $C$ have standard Gaussian law $\gamma$.
To avoid confusion between a random variable and concentration, write
$r(c)=R(2,1,c)$ and $d(c)=30c^4-5c^2+247/8$. Proposition~\ref{prop:ratio}
gives
\begin{equation}\label{eq:r}
 r(c)=-\frac{48c}{d(c)},\qquad
 d(c)=30(c^2-1/12)^2+92/3\geq92/3.
\end{equation}
Consequently $r$ is continuous and odd, vanishes only at zero and has
the sign opposite to $c$. In particular
\[
 |r(c)|\leq\frac{36}{23}|c|.
\]
This proves $r\in L^2(\gamma)$ and integrability of its square and
all products with $L^2$ predictors by Cauchy--Schwarz. Gaussian
symmetry gives $\E_\gamma r=0$. Since a nondegenerate Gaussian has
no atom at zero, $r(C)^2>0$ almost surely, so
\begin{equation}\label{eq:v}
 0<v:=\E_\gamma r(C)^2<\infty.
\end{equation}
The exact integral need not be evaluated to prove this strict
inequality. Likewise $r(c)c=-48c^2/d(c)<0$ off zero, and its
integral is finite and strictly negative.

\begin{theorem}\label{thm:conditional}
If $F:\R\to\mathsf Y$ is any measurable even observation and
$\mathcal A=\sigma(F(C))$, then
\[
 \E_\gamma[r(C)\mid\mathcal A]=0\quad\text{almost surely}.
\]
For every $\mathcal A$-measurable $g\in L^2(\gamma)$,
\begin{equation}\label{eq:error}
 \E_\gamma(r-g)^2=v+\E_\gamma g^2.
\end{equation}
In particular zero is the unique optimal $L^2$ predictor, and its
error is the strictly positive number $v$.
\end{theorem}
\begin{proof}
Every event $A\in\sigma(F)$ is invariant under $c\mapsto-c$:
this holds for preimages of measurable sets and the collection of
symmetric sets is a sigma algebra. Hence $1_A(c)r(c)$ is odd and
integrable, and Gaussian symmetry gives $\int_A r\,d\gamma=0$.
The zero function is $\mathcal A$-measurable and satisfies these
integral identities, so it is a conditional expectation of $r$.
For a bounded $\mathcal A$-measurable $g$, the identities first give
$\E rg=0$ for simple functions and then by bounded approximation.
For general $L^2$ $g$, truncate it and use Cauchy--Schwarz to pass
to the limit. Expanding the square proves~\eqref{eq:error}.
Its extra term is nonnegative, with equality exactly when $g=0$
in $L^2$.
\end{proof}

Our actual observations are
\begin{equation}\label{eq:S}
 E(c)=10+8c^2,\qquad
 H(c)=\frac{2(17+16c^4)}{(10+8c^2)^2},\qquad S(c)=(E(c),H(c)).
\end{equation}
They are continuous and even. Moreover $E$ recovers $|c|$ by
$|c|=\sqrt{(E-10)/8}$ on its range $E\geq10$. Thus
$\sigma(S)=\sigma(E)=\sigma(|C|)$ on this experiment.
This identifies exactly which information is retained, rather than
merely showing that some statistic is even.

\begin{corollary}\label{cor:variance}
Define, for $e\geq10$,
\[
 J(e)=\frac{2304((e-10)/8)}
       {\left(30((e-10)/8)^2-5((e-10)/8)+247/8\right)^2}.
\]
Then $r(c)^2=J(E(c))$ and
\[
 \Var_\gamma(r(C)\mid S)=J(E(C))>0
 \quad\text{almost surely}.
\]
\end{corollary}
\begin{proof}
Substitute $(E(c)-10)/8=c^2$ in~\eqref{eq:r} squared. The square
is integrable and already measurable in $E$, so its conditional
expectation is itself. The conditional mean is zero by
Theorem~\ref{thm:conditional}. Conditional variance is their difference,
and positivity follows from $C\ne0$ almost surely.
\end{proof}

The same parity calculation determines the first nonzero Hermite
index. Let $H_n$ be the probabilists' polynomials with
$H_0=1$, $H_1=c$ and $H_{n+1}=cH_n-nH_{n-1}$. Induction gives
$H_n(-c)=(-1)^nH_n(c)$. Polynomial Gaussian moments and
\eqref{eq:r} imply integrability of $H_nr$. Therefore
\[
 \E_\gamma[H_{2n}(C)r(C)]=0,\qquad
 \E_\gamma[H_1(C)r(C)]=\E_\gamma[Cr(C)]<0.
\]
The mean coefficient is zero and the degree-one coefficient is nonzero:
the Hermite rank is exactly one. This does not contradict invisible
prediction from $S$, which observes only the even coordinate $|C|$.
It is a parity statement about a different observation space.

\section{Identical observed laws with opposite mean transfer}
Put
\[
 b(c)=\frac{c}{1+|c|},\qquad
 d\gamma_\varepsilon(c)=(1+\varepsilon b(c))\,d\gamma(c),
 \qquad |\varepsilon|\leq1.
\]
The use of $b$ here is local to this probability section; the fixed
second triad amplitude remains one. Since $b$ is continuous, odd and
$|b(c)|<1$, the density is strictly positive at every finite $c$,
bounded above by two, and integrates to one. Thus these are genuine
probability laws, equivalent to $\gamma$, including the endpoints
$\varepsilon=\pm1$. All $\gamma$-integrable functions remain
integrable under them.

\begin{theorem}\label{thm:tilts}
For every measurable even map $F$, the full pushforward laws satisfy
$F_\#\gamma_\varepsilon=F_\#\gamma$. In particular
\[
 S_\#\gamma_1=S_\#\gamma_{-1}=S_\#\gamma.
\]
Nevertheless, for the finite positive constant
\[
 m=-\E_\gamma[b(C)r(C)]>0,
\]
one has $\E_{\gamma_\varepsilon}r=-\varepsilon m$.
Thus $\gamma_1$ and $\gamma_{-1}$ are distinct and give opposite
strict mean transfer while their joint observed laws agree exactly.
\end{theorem}
\begin{proof}
For a measurable observation event $D$, its preimage $A=F^{-1}(D)$
is symmetric. The integrable function $1_A b$ is odd, so
\[
 \gamma_\varepsilon(A)=\int_A(1+\varepsilon b)\,d\gamma
                         =\gamma(A).
\]
This is equality on every measurable set of the observation space,
which proves equality of the pushforwards, not just moment matching.
The product $br$ is integrable because $b$ is bounded and $r$ is
integrable. Formula~\eqref{eq:r} gives
\[
 b(c)r(c)=-\frac{48c^2}{(1+|c|)d(c)}<0\qquad(c\ne0).
\]
Its negative is positive almost surely and integrable, so $0<m<\infty$.
Finally integrate $r(1+\varepsilon b)$ and use $\E_\gamma r=0$.
Different means of one integrable observable prove that the two
endpoint laws are distinct.
\end{proof}

The deterministic case already has the same sign reversal at $c=1$
and $c=-1$, with unchanged observed values. The probability statement
is stronger: it preserves the distribution of every even measurable
statistic under strictly positive changes of density, rather than
pairing only two deterministic samples. It excludes determining this
mean transfer from the law of $S$ alone over the displayed family of
experiments. It does not exclude a closure which additionally specifies
a conditional phase law or dynamical assumptions.

\section{Conditional effects of the bias and a common-risk identity}
The following consequences are proved here from the mapped formulas;
they are not attributed to additional named Lean declarations.
They clarify how conditional mean and uncertainty change even when
the observation law is completely unchanged.

For $r_0=|c|$, the products $b(c)r(c)$ and $r(c)^2$ are even and
are measurable in $E(c)$. The tilted conditional expectations are
\begin{equation}\label{eq:biasedconditional}
 \E_{\gamma_\varepsilon}[r\mid S]=\varepsilon br,
 \qquad
 \E_{\gamma_\varepsilon}[r^2\mid S]=r^2.
\end{equation}
To prove the first identity, test against any bounded
$\sigma(S)$-measurable $g$. Such a function is even, and
\[
 \E_{\gamma_\varepsilon}[rg]
   =\E_\gamma[rg]+\varepsilon\E_\gamma[brg]
   =\varepsilon\E_\gamma[brg].
\]
Since $brg$ is even, its integral under $\gamma_\varepsilon$ equals
the one under $\gamma$ by Theorem~\ref{thm:tilts}. The right side
is therefore $\E_{\gamma_\varepsilon}[(\varepsilon br)g]$.
All terms are integrable, giving the defining conditional identities.
For the second formula, $r^2$ is already integrable and observed.
It follows that
\begin{equation}\label{eq:biasedvariance}
 \Var_{\gamma_\varepsilon}(r\mid S)
       =r^2(1-\varepsilon^2b^2)>0\quad\text{almost surely}.
\end{equation}
Here strict positivity also holds at $|\varepsilon|=1$, because
$|b(c)|<1$ for every finite $c$ and $r(c)\ne0$ off zero.
Equivalence of the tilted Gaussian laws preserves that null exception.

There is also a prediction obstruction if one must use the same
observation rule without being told which endpoint law applies.
Let $g=\varphi(S)$ with $\E_\gamma g^2<\infty$, and write
$\mathcal R_\pm(g)=\E_{\gamma_{\pm1}}(r-g)^2$.
Bounded densities ensure these errors are finite. Averaging the two
densities cancels the odd bias exactly, so
\begin{equation}\label{eq:minimax}
 \frac{\mathcal R_+(g)+\mathcal R_-(g)}2
       =\E_\gamma(r-g)^2=v+\E_\gamma g^2\geq v.
\end{equation}
Thus at least one of the two errors is at least $v$. The common rule
$g=0$ has error exactly $v$ under both laws, because $r^2$ is even.
Consequently the minimum over common rules of
$\max\{\mathcal R_+(g),\mathcal R_-(g)\}$ is exactly $v$.
This is an elementary two-experiment consequence, with a complete
proof in~\eqref{eq:minimax}; no new minimax principle is claimed.
The best rule when $\varepsilon$ is known is instead
$\varepsilon br$, as~\eqref{eq:biasedconditional} states. The
common-rule obstruction must not be read as its failure.


\section{The actual retained dynamics and its Gaussian law}\label{sec:retainedflow}
There is a rigorous dynamical companion for this same finite family.
It evolves all three amplitudes rather than holding $a=2,b=1$ fixed.
We now use the state $(a,b,c)\in\R^3$ and the product law
$\Gamma=\gamma\otimes\gamma\otimes\gamma$. This is a different
experiment from the one-parameter biased laws above.

\subsection{Projection and the global solution}
For $k\ne0$ let $\Pi_k v=v-k(k\cdot v)/|k|^2$. The raw vector
convolution corresponding to the negative complex advection under our
$iV$ convention is $iN_k$, where
\[
 N_k=\sum_{p+q=k}(q\cdot V(p))V(q).
\]
The same enumeration used in Lemma~\ref{lem:table} gives
\[
 N_{e_1}=(-2bc,bc,0),\quad N_{e_2}=(ac,-2ac,0),\quad
 N_{e_1+e_2}=(-ab,-ab,0).
\]
For example at $e_1$ the two terms are $(-bc,0,0)$ and
$(-bc,bc,0)$; at $e_2$ they are $(0,-ac,0)$ and $(ac,-ac,0)$.
At $e_1+e_2$ they are $(-ab,0,0)$ and $(0,-ab,0)$.
The opposite outputs have opposite vector values. Applying the
projectors removes the $e_1$ component of the first vector, the
$e_2$ component of the second, and the whole third vector.
The retained coefficients therefore evolve by
\begin{equation}\label{eq:ode}
 \dot a=bc,\qquad \dot b=-ac,\qquad \dot c=0.
\end{equation}
The Fourier phase follows directly from $-(iV(p)\cdot iq)iV(q)
=i(q\cdot V(p))V(q)$. Thus the sign of the displayed drift is
fixed by the actual synthesis convention.

\begin{proposition}\label{prop:flow}
For every real $t$, the global solution map of~\eqref{eq:ode} is
\begin{equation}\label{eq:flow}
 \Phi_t(a,b,c)=\bigl(a\cos(tc)+b\sin(tc),
                  -a\sin(tc)+b\cos(tc),c\bigr).
\end{equation}
It is a smooth group of invertible transformations, preserves
$\alpha(a^2+b^2)+\beta c^2$ for all real $\alpha,\beta$, and
preserves the actual probability law $\Gamma$.
\end{proposition}
\begin{proof}
Differentiate the three coordinates to verify~\eqref{eq:ode} and
the initial value. The polynomial vector field is locally Lipschitz,
so uniqueness identifies this explicit global solution. The addition
formulas give $\Phi_{s+t}=\Phi_s\Phi_t$ and $\Phi_{-t}\Phi_t=I$.
The first two coordinates are an orthogonal rotation at fixed $c$;
their squared sum and $c$ are unchanged, proving every displayed
quadratic invariant. The transformation is nonlinear on the full
state: at $t=\pi/2$, the first coordinate of
$\Phi_t(1,0,1)$ is zero, whereas that of
$\Phi_t(1,0,0)+\Phi_t(0,0,1)$ is one.

For any bounded Borel test $h$, integrate first in $(a,b)$ at fixed
$c$. The two-dimensional standard Gaussian is invariant under every
orthogonal rotation, as follows from its radial density and the
unit absolute Jacobian. Therefore
\[
 \int h(\Phi_t(a,b,c))\,d\gamma(a)d\gamma(b)
      =\int h(a,b,c)\,d\gamma(a)d\gamma(b).
\]
Fubini then integrates over $c$, proving $\Phi_{t\#}\Gamma=\Gamma$.
This proves invariance on all Borel sets by indicator tests, not just
constancy of finitely many Gaussian moments.
\end{proof}

In particular the weighted energy $E=2a^2+2b^2+8c^2$ and the
actual unweighted kinetic norm $2a^2+2b^2+4c^2$ are both conserved
by this retained inviscid flow. Along any nonzero trajectory,
differentiate the concentration in~\eqref{eq:ratio}; because $E$ is
constant, its derivative is
\[
 \frac{d\mathcal C}{dt}
  =\frac{2}{E^2}(4a^3bc-4ab^3c)
  =\frac{8abc(a^2-b^2)}{E^2}=\Sigma.
\]
Consequently the original diagnostic is exactly
$R=-d\log(\mathcal C-1/32)/dt$ \emph{for the retained ODE}.
The positive denominator was already proved. This elementary
derivative identification is a human deduction from the finite table
and actual drift; it is not a claim that the full fluid solution
stays on these six modes. The static slice $a=2,b=1$ is also not
an invariant family of this flow.

\subsection{An original Gaussian gradient obstruction}
Let $G_1,G_2,G_3$ be the first three coordinates on the original
countable Gaussian product space. Let $\mathbf e_j$ be the
corresponding orthonormal source directions, and consider the actual
square-integrable vector field
\[
 U=G_2G_3\mathbf e_1-G_1G_3\mathbf e_2.
\]
Its first three components are precisely the drift~\eqref{eq:ode}
evaluated at these coordinates. Products here are finite Gaussian
polynomials; no unrestricted multiplication of arbitrary $L^2$
functions is assumed. Independence gives
\begin{equation}\label{eq:U-norm}
 \|U\|_2^2=\E(G_2^2G_3^2+G_1^2G_3^2)=2.
\end{equation}

Define $D$ to be the original closed Gaussian derivative obtained by
closing the cylindrical polynomial derivative. On that polynomial
core, one-coordinate Gaussian integration by parts gives
\[
 \E[\partial_1 f\,G_2G_3]=\E[fG_1G_2G_3]
       =\E[\partial_2 f\,G_1G_3].
\]
To justify each integration, fix all remaining finite coordinates;
the Gaussian density makes the polynomial boundary term zero, and
all absolute polynomial moments are finite, so Fubini applies.
Subtracting proves $\langle Df,U\rangle=0$ on the core. If
$f$ lies in the original closed domain, take core $f_j$ with
$f_j\to f$ and $Df_j\to Df$ in $L^2$. Pairing with the fixed
$U\in L^2$ passes the identity to the limit. Thus
\begin{equation}\label{eq:U-orthogonal}
 \langle Df,U\rangle=0,\qquad
 \|U-Df\|_2^2=2+\|Df\|_2^2\quad(f\in\operatorname{Dom}D).
\end{equation}
In particular $U$ is not an original closed gradient, and its exact
distance from the gradient range or its closure is $\sqrt2$,
attained by the zero gradient. Equivalently $U$ belongs to the
adjoint divergence domain with divergence zero: the functional
$f\mapsto\langle Df,U\rangle$ is represented by the zero scalar.
The norm in~\eqref{eq:U-norm} proves this is a nonzero field.

For cylindrical polynomial $f,g$, apply the same identity to $fg$
and use the ordinary product rule:
\[
 \E[f\langle Dg,U\rangle]=-\E[g\langle Df,U\rangle].
\]
These pairings are integrable finite Gaussian polynomials. This is a
skew transport pairing on the polynomial core, not an identification
of the full generator domain of the Koopman group. All assertions
refer to the original closed derivative, not a newly chosen maximal
operator. They reproduce the source's nonzero quadratic second-chaos
field and its closed-gradient obstruction.

\section{Coordinate dephasing without mixing or norm decay}\label{sec:dephasing}
Under the invariant law $\Gamma$, write $A_t$ for the first coordinate
of $\Phi_t(A,B,C)$, where $A,B,C$ are independent standard Gaussians.
By invariance, $A_t$ remains centered of variance one. The actual
two-time correlation, however, is
\begin{equation}\label{eq:dephasing}
 \E[A_tA]=\E[\cos(tC)]\E[A^2]
             +\E[\sin(tC)]\E[AB]=e^{-t^2/2}.
\end{equation}
The products are integrable by Cauchy--Schwarz or bounded sine/cosine
and finite Gaussian second moments, so independence and Fubini
justify the factorization. For completeness let
$\chi(t)=\E e^{itC}$. Dominated differentiation is allowed by
$\E|C|<\infty$. Gaussian integration by parts gives
$\E[Ce^{itC}]=it\chi(t)$, hence $\chi'=-t\chi$ and $\chi(0)=1$.
Solving gives $\chi(t)=e^{-t^2/2}$ and proves the last equality.
Thus the displayed coordinate loses its correlation at a Gaussian rate.

There is nonetheless a centered variance-one conserved observable:
$C\circ\Phi_t=C$. Consequently
\begin{equation}\label{eq:fixedcorrelation}
 \E[(C\circ\Phi_t)C]=1\qquad\text{for all }t.
\end{equation}
This is enough to disprove decay of every centered $L^2$ correlation.
It also excludes mixing, whether defined using all centered $L^2$
tests or using events: the nontrivial event $\{C>0\}$ is invariant
and has probability $1/2$, so its self-intersection has probability
$1/2$ rather than the mixing limit $1/4$.

Finally the composition operators $\mathcal U_t f=f\circ\Phi_t$ on
$L^2(\Gamma)$ are norm preserving by the proved invariance, and
$\mathcal U_{-t}$ is their inverse. Thus they form an orthogonal
group on real $L^2$, or a unitary group after complexification.
They preserve means and cannot satisfy
\[
 \|\mathcal U_t f\|_2\leq q\|f\|_2\quad
 \text{for every centered }f,\qquad q<1.
\]
Taking $f=C$ already contradicts this bound; norm preservation also
contradicts it for every nonzero centered $f$ separately. Decay of
one correlation in~\eqref{eq:dephasing} is compatible with conservation
of the whole $L^2$ norm. The composition-operator mechanism is
classical Koopman theory~\cite{koopman}; the Gaussian rate is the
ordinary Gaussian characteristic function, not a new mixing theorem.


\section{Viscous dissipation without a uniform pair contraction}
\label{sec:viscoustriad}
The same retained system has an exact viscous counterpart. This is
still the projected finite Fourier system, not a claim that the full
fluid evolution remains supported on the six modes. Their Laplacian
weights are $1,1,2$, so the actual retained equations are
\begin{equation}\label{eq:viscousode}
 \dot a=bc-\nu a,\qquad \dot b=-ac-\nu b,\qquad
 \dot c=-2\nu c.
\end{equation}
Fix $\nu>0$ and put
\[
 \rho(t)=e^{-\nu t},\quad \eta(t)=e^{-2\nu t},\quad
 \vartheta(t)=\frac{1-\eta(t)}{2\nu}.
\]
The complete global solution is
\begin{equation}\label{eq:viscousflow}
 \Phi_t^\nu(a,b,c)=
 \bigl(\rho(t)[a\cos(\vartheta(t)c)+b\sin(\vartheta(t)c)],
       \rho(t)[-a\sin(\vartheta(t)c)+b\cos(\vartheta(t)c)],
       \eta(t)c\bigr).
\end{equation}
Indeed $\rho'=-\nu\rho$, $\eta'=-2\nu\eta$ and
$\vartheta'=\eta$. Differentiating the first rotated coordinate
gives its damping plus the second coordinate multiplied by the
current $c(t)=\eta(t)c$; the second gives damping minus the first
times $c(t)$. The last coordinate solves its scalar equation and
the initial map is the identity. Local uniqueness for the polynomial
vector field identifies this global formula with the solution.
Also
$\vartheta(s+t)=\vartheta(t)+\eta(t)\vartheta(s)$ and the
exponential multiplication laws give
$\Phi_{s+t}^\nu=\Phi_s^\nu\Phi_t^\nu$ directly.

\subsection*{The actual physical metric}
The relevant energy here is the kinetic one
\[
 \mathcal K(a,b,c)=2a^2+2b^2+4c^2,
\]
rather than the enstrophy-weighted $\mathcal E$ in Section 2.
It equals the actual normalized torus integral
$(2\pi)^{-3}\int_{[0,2\pi]^3}|u_{a,b,c}|^2$.
To see this without an assumed Parseval identity, use the displayed
velocity~\eqref{eq:velocity}. The two pure sine basis vectors have
mean square two each; the third has two components and mean square
four. Their cross integrals vanish: integrate
$\sin x_1\sin(x_1+x_2)$ first in $x_2$, and
$\sin x_2\sin(x_1+x_2)$ first in $x_1$; the remaining pair is
pointwise orthogonal. All integrals concern continuous functions on
a compact box, so ordinary Fubini applies. Bilinearity therefore gives
the full physical Gram form
\[
 \langle u_x,u_y\rangle_{L^2(\mathbb T^3)}
 =2x_1y_1+2x_2y_2+4x_3y_3.
\]
The coordinates in this display are numbered $1,2,3$.
Linearity of the velocity family proves that the squared physical
distance is exactly $\mathcal K(x-y)$.

Rotation preserves the squared length of the first two coordinates.
Consequently~\eqref{eq:viscousflow} gives the exact energy identity
\begin{equation}\label{eq:viscousenergy}
 \mathcal K(\Phi_t^\nu(a,b,c))
 =2e^{-2\nu t}(a^2+b^2)+4e^{-4\nu t}c^2
 \leq e^{-2\nu t}\mathcal K(a,b,c)\qquad(t\geq0).
\end{equation}
Every nonzero initial state has strictly smaller energy at every
$t>0$. This decay toward zero does not imply contraction of
distances between arbitrary pairs.

\subsection*{Equal third amplitudes and an expanding pair}
If $x$ and $y$ have the same third amplitude, their first two
coordinates rotate by the same angle, and their evolved third
coordinates remain equal. Thus the exact slice identity is
\[
 \mathcal K(\Phi_t^\nu x-\Phi_t^\nu y)
       =e^{-2\nu t}\mathcal K(x-y).
\]
By contrast, fix any $\nu>0$ and $t>0$. Then $\vartheta(t)>0$.
Set $d=\pi/\vartheta(t)$ and consider
$x=(A,0,0)$, $y=(A,0,d)$. They initially differ only in the
third amplitude. Their evolved first coordinates have opposite signs:
\[
 \Phi_t^\nu x=(\rho A,0,0),\qquad
 \Phi_t^\nu y=(-\rho A,0,\eta d).
\]
Here $\rho=\rho(t)$ and $\eta=\eta(t)$. Therefore
\begin{equation}\label{eq:triadamplification}
 \mathcal K(x-y)=4d^2,\qquad
 \mathcal K(\Phi_t^\nu x-\Phi_t^\nu y)
                  =8\rho^2A^2+4\eta^2d^2.
\end{equation}
For any proposed finite real squared-distance factor $L$, choose
$A=d(|L|+1)/\rho$. Dividing the two positive energies gives
$2(|L|+1)^2+\eta^2>L$. No finite factor controls all initial
pairs at this fixed positive time. In particular no global Lipschitz
constant in the physical kinetic norm exists: a finite nonnegative
norm constant would give a squared-distance factor by squaring.
Taking $L=1$ supplies two states whose separate energies both strictly
decrease by~\eqref{eq:viscousenergy}, while their mutual physical
distance increases. These examples use unrestricted initial amplitudes;
they do not refute estimates on a bounded state set or local smooth
dependence on initial data.

\subsection*{The inviscid Gaussian invariance also changes}
The product standard Gaussian $\Gamma$ from Section~\ref{sec:retainedflow} is not invariant
under this deterministic viscous flow. If the initial third coordinate
$C$ is standard normal, its evolved second moment is
\[
 \E[(\eta(t)C)^2]=e^{-4\nu t}<1\qquad(t>0).
\]
Thus equality of the pushforward law with $\Gamma$ is impossible.
Equivalently the initially centered $L^2(\Gamma)$ observable
$h(a,b,c)=c^2-1$ has evolved mean $e^{-4\nu t}-1<0$.
These expectations are legitimate finite Gaussian moments; the first
two components in~\eqref{eq:viscousflow} are bounded in absolute
value by $\rho(t)(|a|+|b|)$, so no missing integrability assumption
is concealed. There is no noise or compensating forcing in
\eqref{eq:viscousode}. The earlier invariant-Gaussian and unitary
composition statements concern the inviscid flow and remain distinct.

The original declarations are in
\par\noindent\code{ViscousTriadFlow},
\code{ViscousTriadSeparation} and \code{TriadPhysicalPairing}.
They supply the actual finite ODE, exact energy identities, physical
torus metric, and unrestricted-amplitude obstruction. None of these
statements is a continuum Navier--Stokes stability, invariant-measure
or regularity theorem.

\section{Primary comparison and formal correspondence}
Fourier triad transfer and phase dependence are established topics.
Moffatt~\cite{moffatt} was inspected in the primary PDF through the
introduction and the two-dimensional representation in Section 2;
its distinction between instantaneous triad data and exact evolution
is retained here. Murray--Bustamante~\cite{murray} was checked at
primary abstract level for phase dynamics with amplitudes held fixed
in a Burgers model. Benavides--Bustamante~\cite{benavides}, submitted
May 2026, was checked at primary abstract level for a stochastic phase
closure in two-dimensional turbulence. Its statistical and dynamical
assumptions are not hypotheses of our static experiment. Our
nonidentifiability example is not presented as a counterexample to
either paper or to a fully specified turbulent closure.

Parity orthogonality, conditional expectation as $L^2$ projection and
bounded changes of density are classical. The contribution is a
fully reproducible exact family combining the actual partition,
solenoidal representative, signed transfer denominator and full-law
information loss. Priority of the particular formulas and combined
formalization remains unresolved after this bounded comparison. Koopman's primary 1931 pages 315--316 were inspected for invariant-measure composition and unitary action. The rotational flow, skew Gaussian integration by parts and Gaussian characteristic function are classical mechanisms; no new general principle of dephasing or Hodge decomposition is claimed.

All formal sources are frozen at commit
\begin{center}\small\ttfamily
bab688bb6140cd6316c80680095a914ea0d346db.
\end{center}
The accompanying theorem map records exact declaration names,
headers, lines, file hashes and permanent links. The main modules are:
\begin{longtable}{@{}p{0.28\textwidth}p{0.66\textwidth}@{}}
Actual finite family&\code{ParameterizedSectorTriad}\\[3pt]
Directional cells&\code{AngularTriadPartition}\\[3pt]
Real representative&\code{ParameterizedTriadRepresentative}\\[3pt]
Gaussian rate and signs&\code{GaussianTriadPhase}\\[3pt]
Hermite parity/rank&\code{GaussianTriadHermite}\\[3pt]
Positive biased laws&\code{GaussianTriadBias}\\[3pt]
Conditional mean/variance&\code{GaussianTriadConditional}\\[3pt]
Retained drift and Gaussian field&\code{DictionaryTriadPotential}, \code{QuadraticTriadGaussian}\\[3pt]
Global flow and dephasing&\code{TriadGaussianFlow}, \code{TriadGaussianDephasing}\\
\end{longtable}
The map distinguishes the original formal statements from the
classical sigma-field, endpoint-density and biased-conditional/common-risk
deductions in Sections 4--6. The derivative identity along the retained
ODE is also proved explicitly. The source conditional-variance theorem
concerns the original symmetric Gaussian, and is not misattributed to
the tilted laws.

Original build/axiom evidence and source hashes were checked at the
frozen baseline; no new Lean build or mathematical source change
is claimed. Exact finite algebra checks independently reconstruct
the triad table, labels, denominator and scalar identities. They do
not replace the printed measure-theoretic proofs with Gaussian
sampling or numerical quadrature.

The invariant product Gaussian belongs to the finite retained ODE.
No invariant full-fluid measure, turbulent closure, continuum regularity
or singularity conclusion is inferred.

\begin{thebibliography}{9}
\bibitem{koopman} B. O. Koopman, Hamiltonian systems and transformations
in Hilbert space, \emph{Proc. Natl. Acad. Sci. USA} \textbf{17} (1931), 315--318.
\url{https://new.math.uiuc.edu/MathMLseminar/seminarPapers/Koopman1931.pdf}.

\bibitem{moffatt} H. K. Moffatt, Note on the triad interactions of
homogeneous turbulence, \emph{J. Fluid Mech.} \textbf{741} (2014), R3.
\href{https://doi.org/10.1017/jfm.2013.637}{doi:10.1017/jfm.2013.637}.
Primary PDF: \url{https://www.damtp.cam.ac.uk/user/hkm2/PDFs/Moffatt2014a.pdf}.
\bibitem{murray} B. P. Murray and M. D. Bustamante, Energy flux
enhancement, intermittency and turbulence via Fourier triad phase
dynamics in 1D Burgers equation, primary preprint arXiv:1705.08960v2
(2018). \url{https://arxiv.org/abs/1705.08960}.
Abstract-level comparison only.
\bibitem{benavides} S. J. Benavides and M. D. Bustamante,
Triad phase dynamics determine cascade direction in two-dimensional
turbulence, primary preprint arXiv:2605.03049v1 (2026).
\url{https://arxiv.org/abs/2605.03049}. Abstract-level comparison only.
\end{thebibliography}
\end{document}
